\section{Related Work}
\label{sec:related-work}

\paragraph{Learned computation and continual organization.}
The Transformer couples learned projections, attention-based information
selection, and feedforward transformations to build compositional sequence
models~\cite{vaswani2017attention}. Prometheus shares the use of neural
representations and task-driven credit assignment; its proposed distinction
concerns how executable operations change and remain reusable during later
learning. Continual learning already offers architectural mechanisms for this
problem. DEN selectively retrains relevant subnetworks, conditionally expands
capacity, and uses splitting and timestamped inference to manage semantic
drift~\cite{yoon2018lifelong}. Thus growth, sharing, and retention are established
capabilities. A comparison with such methods must identify the additional work
performed by loop semantics and include each method's stored capacity and
maintenance costs.

\paragraph{Learning through discrete changes.}
Stochastic computation graphs combine differentiable computation with
score-function gradient terms and support surrogate-loss
implementations~\cite{schulman2015stochastic}. For discontinuous models, Lee
et al. explicitly separate regionwise reparameterization derivatives from
boundary contributions~\cite{lee2018reparameterization}. Our coupled
within-class and boundary update uses this established distinction: ordinary
neural credit learns encoders, readouts, and calls, while class-boundary
feedback changes the distribution over executable loop actions. The
derivative decomposition itself is not a novelty claim. Its value here must
come from the computational organization and a measured advantage over
equally informed ordinary boundary optimization or search.

\paragraph{Topology and reversible computation.}
Topology has an established computational role. Kitaev's anyon construction
uses braiding to implement transformations in a physically protected quantum
system~\cite{kitaev2003fault}. Our present algebraic-cover construction instead
uses continuation to define inspectable state permutations; digital simulation
does not inherit physical fault tolerance or energy savings. The relevant
question is whether homotopy-preserving implementation freedom helps a
learning system acquire and retain useful computations.

Permutation operations also delimit the current computational domain.
Bennett's reversible simulation retains auxiliary information, copies the
result, and reverses the computation to clean its history~\cite{bennett1973logical}.
It provides a constructive route toward irreversible tasks while making the
extra workspace and computation explicit. Such an extension would need to
preserve the same learning interface and account for those resources;
reversible representability alone does not establish a general trainable
Block or an efficient acquisition mechanism.
